\documentclass[11pt]{article}
\usepackage{mathblatt}
% Präambel des Setzers (werkzeuge/setzer-praeambel.tex, wortgleich); dazu \gteil (Muster T2U)
\makeatletter
\setlength{\mbpfbe}{0mm}% keine Punkte (pruefung.md § 5: nur beim Prüfstein)
\newlength{\gEin}\setlength{\gEin}{0pt}
\let\gSchrift\relax
\renewcommand{\pfkreuzzeile}[1]{\par\noindent\hangindent3.2em\hangafter1\hspace*{1.6em}$\square$\ #1\par}
\newcommand{\gkreuz}[1]{$\square$\ #1\hspace{2em}}
\newcommand{\gkreuzzeile}[2]{\par\noindent\hangindent2.2em\hangafter1\mbox{#1\,$\square$}\ #2\par}% Buchstabe, Option: eine je Zeile (neben Skizze)
% links, was man ansieht, rechts, was man tut; Spaltengrenze überall an derselben Stelle (gemeinsam.md „Satz“)
\newsavebox{\gbox}\newlength{\gLinks}
% v0.10 (Lehrer 13.10.): keine starre Halbierung – die linke Spalte ist so breit wie die Figur (+ 6 mm), die
% Antwort steht direkt daneben; ist die Figur breiter als 0,62 Zeile, steht die Antwort darunter
\newcommand{\gzwei}[2]{\par\smallskip\noindent\sbox{\gbox}{#1}%
  \setlength{\gLinks}{\dimexpr\wd\gbox+6mm\relax}%
  \ifdim\gLinks>0.62\linewidth
    \usebox{\gbox}\par\smallskip\noindent\begin{minipage}{\linewidth}#2\end{minipage}%
  \else
    \begin{minipage}[t]{\gLinks}\vspace{0pt}\usebox{\gbox}\end{minipage}%
    \begin{minipage}[t]{\dimexpr\linewidth-\gLinks\relax}\vspace{0pt}\raggedright #2\end{minipage}%
  \fi\par}
% v0.9 (12.10., pruefung.md § 5 „Marke mit B, Bedingungen rechts“): drittes Argument von pfaufg = Bedingungen
% der Prüfung (v0.10: \pfohnetr = TR durchgestrichen, „FOR“), klein grau im rechten Rand auf der Grundlinie der Nummer, Zeilen mit \\
\newcommand{\gbedrand}[2]{\ifblank{#2}{}{\rlap{\hspace*{\dimexpr#1+2mm\relax}%
  \parbox[t]{15mm}{\raggedright\scriptsize\leavevmode\color{mbgrau}\strut#2}}}}
% Teil eines Bündels (a, b …): Marke links im Rand, Bedingungen rechts im Rand, auf der Zeile des Teils
\newcommand{\gteilrand}[2]{\llap{\mbpfrandmarke{#1}\hspace*{\mbpfnr}}\gbedrand{\linewidth}{#2}}
% Teil eines Bündels: Buchstabe in eigener Spalte, Inhalt eingerückt (\gEin = Einzug, damit die Spaltengrenze von
% \gzwei an derselben Stelle bleibt)
\newenvironment{gbteil}[3]{\par\smallskip\noindent\gteilrand{#2}{#3}\makebox[\mbpfteil][l]{\textbf{#1}}%
  \begin{minipage}[t]{\dimexpr\linewidth-\mbpfteil\relax}\setlength{\gEin}{\mbpfteil}\raggedright\strut\ignorespaces}%
 {\end{minipage}\par}
% Grafik höchstens so breit wie die Spalte (Bündel nebeneinander, Koordinatensystem ohne Karo)
\newsavebox{\gpbox}
\newcommand{\gpassend}[1]{\sbox{\gpbox}{#1}\ifdim\wd\gpbox>\linewidth\typeout{GPASSEND: verkleinert \the\wd\gpbox\space auf \the\linewidth}\resizebox{\linewidth}{!}{\usebox{\gpbox}}\else\usebox{\gpbox}\fi}
% Skizzenfeld (v0.10, Lehrer 13.10.): Karo ab dem linken Rand über die Zeile, „Skizze“ klein grau oben links im Feld
% (wie \pfskizzenfeld der Vorlage); der rechte Rand bleibt den Bedingungen
\newcommand{\gskizze}[1]{\par\smallskip\noindent\begin{tikzpicture}
  \clip (0,0) rectangle ({\linewidth-1mm},{#1*10mm});
  \draw[black!16,line width=0.3pt,step=5mm] (0,0) grid ({\linewidth},{#1*10mm});
  \node[anchor=north west,inner sep=2pt,font=\scriptsize,text=mbgrau,fill=white] at (0,{#1*10mm}) {Skizze};
  \end{tikzpicture}\par}
\RenewDocumentEnvironment{pfaufg}{m m m}{%
  \begin{lrbox}{\mbaufgabebox}\begin{minipage}[t]{\linewidth}%
  \gSchrift\noindent\hspace*{\gEin}\mbpfrandmarke{#1}\gbedrand{\linewidth-\gEin}{#3}%
  \begin{minipage}[t]{\mbpfnr}\strut\textbf{#2}\end{minipage}%
  \begin{minipage}[t]{\dimexpr\linewidth-\gEin-\mbpfnr\relax}\raggedright\strut\ignorespaces}%
 {\end{minipage}%
  \end{minipage}\end{lrbox}%
  \setlength{\mbaufgabehoehe}{\dimexpr\ht\mbaufgabebox+\dp\mbaufgabebox\relax}%
  \ifdefempty{\gkopfnach}{\par\addvspace{5pt}\Needspace*{\mbaufgabehoehe}}{%
    % Stufenkopf, der mit dieser (ersten) Aufgabe zusammenbleibt: Platz für beide verlangen, dann den Kopf setzen
    % (12.10., Quadratische: ein Kopf stand allein unten, die Aufgabe mit Koordinatensystem auf der nächsten Seite)
    \par\Needspace*{\dimexpr\mbaufgabehoehe+40pt\relax}\gkopfnach\gdef\gkopfnach{}\par\addvspace{5pt}}%
  \noindent\usebox{\mbaufgabebox}\par\addvspace{5pt}}
% Stufenkopf vormerken: \gkopfvor{<Kopf>} – die nächste pfaufg setzt ihn, sobald beide zusammen Platz haben
\def\gkopfnach{}
\newcommand{\gkopfvor}[1]{\gdef\gkopfnach{#1}}
% Teilaufgabe mit Skizze links in gzwei: Buchstabe in eigener Spalte, darunter Text und Skizze
\newcommand{\gteil}[2]{\makebox[\mbpfteil][l]{\textbf{#1}}\parbox[t]{\dimexpr0.5\textwidth-\gEin-\mbpfnr-5mm-\mbpfteil\relax}{\raggedright #2}}
% Stufenkopf der Vorlage (mathblatt.sty 2026-10-12: Linie, Abstand, 13 pt, \pfjahre) ohne ihren Vorhalt von
% 14 Zeilen: der Setzer fordert vor dem Kopf 0,16 Seitenhöhe (Kopf + erste Aufgabe), mehr nicht
\setlength{\mbpfstufevorhalt}{0pt}
\makeatother
\newcommand{\karo}[1]{\par\smallskip\noindent\begin{tikzpicture}
  \clip (0,0) rectangle ({\linewidth-1mm},{#1*10mm});
  \draw[black!16,line width=0.3pt,step=5mm] (0,0) grid ({\linewidth},{#1*10mm});
  \end{tikzpicture}\par}
\newcommand{\kk}[1]{$\square$\,#1\hspace{0.9em}}
\newcommand{\linie}[1]{\,{\color{black!45}\rule[-1pt]{#1}{0.4pt}}\,}
\newcommand{\pkt}{\textperiodcentered{}}
% v0.14 (14.10.): Stufenkopf mit fünf Punkten (Option = Jahre für die Bestellkarte) oder „zuletzt 19“
\NewDocumentCommand{\gpunkte}{o m}{#2}
\newcommand{\gzuletzt}[1]{zuletzt #1}
% v0.7 (12.10.): Schülerfrage unter dem Stufenkopf; Ankreuzoption als Bild (Buchstabe und Kästchen darunter)
\newcommand{\gfrage}[1]{\par\nobreak\noindent\hspace*{\mbpfnr}{\small\itshape #1}\par\nobreak\addvspace{2pt}}
\newlength{\gbildbreite}\setlength{\gbildbreite}{0.32\linewidth}% Setzer setzt sie je Reihe (vier Bilder: 0,24)
\newcommand{\gbildoption}[3]{\begin{minipage}[t][#3][b]{\gbildbreite}\centering #2\par\smallskip\mbox{#1\,$\square$}\end{minipage}}
% Lineare (12.10.): Bild mit Buchstabe ohne Kästchen (Diagramm zuordnen); Aussage mit „wahr / falsch“ zum Ankreuzen
\newcommand{\gbildname}[3]{\begin{minipage}[t][#3][b]{\gbildbreite}\centering #2\par\smallskip\mbox{\textbf{#1}}\end{minipage}}
\newcommand{\gwahrfalsch}[2]{\par\noindent\hangindent2.2em\hangafter1\makebox[2.2em][l]{(#1)}#2\hfill\mbox{wahr $\square$\quad falsch $\square$}\par}
\begin{document}

\pfheftstil
\pflaufkopf[Abitur GK \pkt{} Prüfung]{Geraden und Ebenen}{AG40}\pfseitevon
\pfheftkopf{Geraden und Ebenen}{}
\pfstufenzeile{Punktprobe an der Geraden \pkt{} Ebenengleichung gesucht \pkt{} Punktprobe an der Ebene}

% ===================== Stufe: Punktprobe an der Geraden =====================
\gkopfvor{\pfstufekopf{Punktprobe an der Geraden}{\gpunkte[25 23]{\textcolor{black!25}{$\circ$}\,\textcolor{black!80}{$\bullet$}\,\textcolor{black!25}{$\circ$}\,\textcolor{black!80}{$\bullet$}\,\textcolor{black!25}{$\circ$}}}\special{pdf:dest (stufe-punktprobe-an-der-geraden) [@thispage /XYZ @xpos @ypos null]}}
% kein Einstieg: keine Schülerfrage in der Gliederung (Schülerfrage, -Fälle, -Form, -Weg, -Skizze; oder Frage) – Einstieg nur von dort, nie aus der Bank (pruefung.md § 5, 12.10.)
% 2023-bebb-gk-A1.4a-m1 (Original 2023-bebb-gk-A1.4a, Schwere I, rang 1; herausgelöste Fassung: keine)
\begin{pfaufg}{GK~’23\\A1.4a}{1.}{\pfohnetr}
Gegeben ist die Gerade $g\colon \vec{x} = \begin{pmatrix}1\\4\\-3\end{pmatrix} + s \cdot \begin{pmatrix}2\\0\\3\end{pmatrix}$, $s \in \mathbb{R}$, und der Punkt $A(5 \mid 0 \mid 0)$. Begründe, dass $A$ nicht auf $g$ liegt.
\karo{1}
\fusshilfe{\mbox{1:~nein}}
\end{pfaufg}
% 2025-bebb-gk-A1.2a-m1 (Original 2025-bebb-gk-A1.2a, Schwere I, rang 2; herausgelöste Fassung: keine)
\begin{pfaufg}{GK~’25\\A1.2a}{2.}{\pfohnetr}
Gegeben ist die Gerade $g\colon \vec{x} = \begin{pmatrix}6\\2\\-2\end{pmatrix} + s \cdot \begin{pmatrix}-3\\0\\2\end{pmatrix}$, $s \in \mathbb{R}$, und der Punkt $P(3 \mid 2 \mid 2)$. Zeige, dass $P$ nicht auf $g$ liegt. Gib die Koordinaten eines Punkts $Q$ auf $g$ an, der sich nur in einer Koordinate von $P$ unterscheidet.
\karo{2}
\fusshilfe{\mbox{2:~P liegt nicht auf g; Q(3|2|0)}}
\end{pfaufg}

% ===================== Stufe: Ebenengleichung gesucht =====================
\gkopfvor{\pfstufekopf{Ebenengleichung gesucht}{\gpunkte[26 25 24 22]{\textcolor{black!80}{$\bullet$}\,\textcolor{black!80}{$\bullet$}\,\textcolor{black!80}{$\bullet$}\,\textcolor{black!25}{$\circ$}\,\textcolor{black!80}{$\bullet$}}}\special{pdf:dest (stufe-ebenengleichung-gesucht) [@thispage /XYZ @xpos @ypos null]}}
% kein Einstieg: keine Schülerfrage in der Gliederung (Schülerfrage, -Fälle, -Form, -Weg, -Skizze; oder Frage) – Einstieg nur von dort, nie aus der Bank (pruefung.md § 5, 12.10.)
% Teile eines Originals (2022-bebb-gk-B3): 2022-bebb-gk-B3b-m1, 2022-bebb-gk-B3c-m2
\begin{pfaufg}{}{3.}{}
Ein Kirchturmdach steht auf dem Quadrat ABCD mit $A(0 \mid 0 \mid 0)$, $B(6 \mid 0 \mid 0)$, $C(6 \mid 6 \mid 0)$, $D(0 \mid 6 \mid 0)$. Darüber liegen $E(3 \mid 0 \mid 5)$, $F(6 \mid 3 \mid 5)$, $G(3 \mid 6 \mid 5)$, $H(0 \mid 3 \mid 5)$ und die Spitze $S(3 \mid 3 \mid 10)$; 1 LE entspricht 1 m. Das Dach hat vier gleiche viereckige Dachflächen wie CGSF und vier dreieckige Giebelflächen. Die Ebene L enthält C, G und F.
% 2022-bebb-gk-B3b-m1 (Original 2022-bebb-gk-B3b, Schwere I, rang 3; herausgelöste Fassung: keine)
\begin{gbteil}{a)}{\raisebox{\dimexpr-\height+\ht\strutbox\relax}{\shortstack[r]{GK~’22\\B3b}}\hspace*{\mbpfnr}}{}
Gib eine Gleichung von L in Parameterform an. Zeige, dass S in L liegt.
\karo{2}
\end{gbteil}
% 2022-bebb-gk-B3c-m2 (Original 2022-bebb-gk-B3c, Schwere I, rang 11; herausgelöste Fassung: keine)
\begin{gbteil}{b)}{\raisebox{\dimexpr-\height+\ht\strutbox\relax}{\shortstack[r]{GK~’22\\B3c}}\hspace*{\mbpfnr}}{}
Die Ebene N ist parallel zu L und enthält B. Der Vektor $\left(\begin{smallmatrix}5\\5\\3\end{smallmatrix}\right)$ ist ein Normalenvektor von L. Bestimme eine Koordinatengleichung von N.
\karo{2}
\end{gbteil}
\fusshilfe{\mbox{3:~a)~$\left(\begin{smallmatrix}6\\6\\0\end{smallmatrix}\right)$ + s · $\left(\begin{smallmatrix}-3\\0\\5\end{smallmatrix}\right)$ + t · $\left(\begin{smallmatrix}0\\-3\\5\end{smallmatrix}\right)$; S liegt in L; b)~30}}
\end{pfaufg}
% 2024-bebb-gk-B3d-m1 (Original 2024-bebb-gk-B3d, Schwere II, rang 4; herausgelöste Fassung: keine)
\begin{pfaufg}{GK~’24\\B3d}{4.}{}
Gegeben ist die Pyramide ABCDS mit $A(0 \mid 0 \mid 0)$, $B(3 \mid 0 \mid 0)$, $C(3 \mid 3 \mid 0)$, $D(0 \mid 6 \mid 0)$ und $S(0 \mid 0 \mid 5)$. Die Seitenfläche CDS liegt in der Ebene E. Bestimme eine Gleichung von E in Koordinatenform. (Zur Kontrolle: $E\colon 5x + 5y + 6z = 30$)
\karo{2}
\fusshilfe{\mbox{4:~30}}
\end{pfaufg}
% 2025-bebb-gk-B3b-m1 (Original 2025-bebb-gk-B3b, Schwere II, rang 5; herausgelöste Fassung: keine)
\begin{pfaufg}{GK~’25\\B3b}{5.}{}
Gegeben ist die Pyramide ABCDS mit $A(0 \mid 0 \mid 0)$, $B(3 \mid 3 \mid 0)$, $C(0 \mid 8 \mid 0)$, $D(-3 \mid 3 \mid 0)$ und $S(0 \mid 0 \mid 6)$; ihre Grundfläche ABCD ist ein Drachenviereck. Die Seitenfläche BCS liegt in der Ebene E. Ermittle eine Koordinatengleichung von E. (Zur Kontrolle: $E\colon 5x + 3y + 4z = 24$)
\karo{2}
\fusshilfe{\mbox{5:~24}}
\end{pfaufg}
% 2026-bb-gk-B3e-m1 (Original 2026-bb-gk-B3e, Schwere III, rang 6; herausgelöste Fassung: keine)
\begin{pfaufg}{GK~’26\\B3e}{6.}{}
Gegeben ist die Pyramide ABCDS mit $A(0 \mid -3 \mid 0)$, $B(5 \mid -1{,}5 \mid 0)$, $C(5 \mid 1{,}5 \mid 0)$, $D(0 \mid 3 \mid 0)$ und $S(1 \mid 0 \mid 3)$; sie ist symmetrisch zur xz-Ebene. Die Ebene $z = 1{,}5$ schneidet die Seitenkanten in T, U, V und $W(0{,}5 \mid 1{,}5 \mid 1{,}5)$. Die Trapeze TUVW und BUVC liegen symmetrisch zueinander bezüglich einer Ebene H. Bestimme eine Gleichung von H.
\gzwei{\begin{ksys3}[xyz,x1min=0,x1max=5,x2min=-3,x2max=3,x3min=0,x3max=3,karo=0.3]\rstrecke{0,-3,0}{5,-1.5,0}\rstrecke{5,-1.5,0}{5,1.5,0}\rstrecke{5,1.5,0}{0,3,0}\rstrecke[dashed]{0,3,0}{0,-3,0}\rstrecke{0,-3,0}{1,0,3}\rstrecke{5,-1.5,0}{1,0,3}\rstrecke{5,1.5,0}{1,0,3}\rstrecke{0,3,0}{1,0,3}\rvieleck{0.5,-1.5,1.5;3,-0.75,1.5;3,0.75,1.5;0.5,1.5,1.5}\rpunkt*[above left]{0}{-3}{0}{A}\rpunkt*[below left]{5}{-1.5}{0}{B}\rpunkt*[below right]{5}{1.5}{0}{C}\rpunkt*[above right]{0}{3}{0}{D}\rpunkt*[above left]{1}{0}{3}{S}\rpunkt*[left]{0.5}{-1.5}{1.5}{T}\rpunkt*[below]{3}{-0.75}{1.5}{U}\rpunkt*[right]{3}{0.75}{1.5}{V}\rpunkt*[above]{0.5}{1.5}{1.5}{W}\end{ksys3}}{\karo{2}}
\fusshilfe{\mbox{6:~7,5}}
\end{pfaufg}

% ===================== Stufe: Punktprobe an der Ebene =====================
\gkopfvor{\pfstufekopf{Punktprobe an der Ebene}{\gpunkte[26 25 23 22]{\textcolor{black!80}{$\bullet$}\,\textcolor{black!80}{$\bullet$}\,\textcolor{black!25}{$\circ$}\,\textcolor{black!80}{$\bullet$}\,\textcolor{black!80}{$\bullet$}}}\special{pdf:dest (stufe-punktprobe-an-der-ebene) [@thispage /XYZ @xpos @ypos null]}}
% kein Einstieg: keine Schülerfrage in der Gliederung (Schülerfrage, -Fälle, -Form, -Weg, -Skizze; oder Frage) – Einstieg nur von dort, nie aus der Bank (pruefung.md § 5, 12.10.)
% 2022-bebb-gk-A1.4a-m1 (Original 2022-bebb-gk-A1.4a, Schwere I, rang 7; herausgelöste Fassung: keine)
\begin{pfaufg}{GK~’22\\A1.4a}{7.}{\pfohnetr}
Gegeben ist die Ebene $E\colon 2x - 5y = 0$. Prüfe, ob der Punkt $(5 \mid 2 \mid 3)$ in $E$ liegt.
\karo{1}
\fusshilfe{\mbox{7:~ja}}
\end{pfaufg}
% 2025-bebb-gk-A1.5b-m1 (Original 2025-bebb-gk-A1.5b, Schwere I, rang 8; herausgelöste Fassung: keine)
\begin{pfaufg}{GK~’25\\A1.5b}{8.}{\pfohnetr}
Gegeben sind $P(1 \mid -3 \mid 0)$, $Q(3 \mid 1 \mid 4)$ und die Ebene $E\colon x + 2y + 2z = 13$. Weise nach, dass $Q$ in $E$ liegt und dass $\overrightarrow{PQ}$ ein Normalenvektor von $E$ ist.
\karo{2}
\fusshilfe{\mbox{8:~Q liegt in E; kollinear zum Normalenvektor}}
\end{pfaufg}
% 2026-bb-gk-A1.2a-m1 (Original 2026-bb-gk-A1.2a, Schwere II, rang 9; herausgelöste Fassung: keine)
\begin{pfaufg}{GK~’26\\A1.2a}{9.}{\pfohnetr}
Gegeben sind die Ebene $E\colon 2x - y + 2z = 3$ und die Gerade $g$ durch $P(1 \mid 3 \mid 2)$ und $Q(7 \mid 0 \mid 8)$. Zeige, dass $P$ in $E$ liegt. Begründe, dass $g$ senkrecht auf $E$ steht.
\karo{2}
\fusshilfe{\mbox{9:~P liegt in E; g steht senkrecht auf E}}
\end{pfaufg}
% 2023-bebb-gk-B3f-m1 (Original 2023-bebb-gk-B3f, Schwere III, rang 10; herausgelöste Fassung: keine)
\begin{pfaufg}{GK~’23\\B3f}{10.}{}
Die Punkte $A(6 \mid 0 \mid 0)$, $B(0 \mid 6 \mid 0)$ und $C(0 \mid 0 \mid 6)$ liegen in der Ebene $L_1\colon x + y + z = 6$, die Punkte $D(3{,}5 \mid 0 \mid 0)$, $E(0 \mid 3{,}5 \mid 0)$ und $F(0 \mid 0 \mid 3{,}5)$ in der Ebene $L_2\colon 2x + 2y + 2z = 7$. Sie sind die Ecken des Körpers ABCDEF. Seine Oberfläche besteht aus den Dreiecken ABC und DEF und drei Trapezen in den Koordinatenebenen. Begründe, dass $T(1 \mid 1 \mid 2)$ im Inneren des Körpers liegt.
\gzwei{\begin{ksys3}[xyz,x1min=0,x1max=6,x2min=0,x2max=6,x3min=0,x3max=6,karo=0.2]\rvieleck{6,0,0;0,6,0;0,0,6}\rstrecke{6,0,0}{3.5,0,0}\rstrecke{0,6,0}{0,3.5,0}\rstrecke{0,0,6}{0,0,3.5}\rstrecke[dashed]{3.5,0,0}{0,3.5,0}\rstrecke[dashed]{0,3.5,0}{0,0,3.5}\rstrecke[dashed]{0,0,3.5}{3.5,0,0}\rpunkt*[left]{6}{0}{0}{A}\rpunkt*[below]{0}{6}{0}{B}\rpunkt*[right]{0}{0}{6}{C}\rpunkt*[below left]{3.5}{0}{0}{D}\rpunkt*[below]{0}{3.5}{0}{E}\rpunkt*[left]{0}{0}{3.5}{F}\end{ksys3}}{\karo{2}}
\fusshilfe{\mbox{10:~T diesseits von L$_1$; T liegt zwischen den Ebenen, im Inneren}}
\end{pfaufg}

\pfblattende{10}{26 & 25 & 24 & 23 & 22 & 21 & 20 & 19 & 18 & 17}{A1.2a B3e & A1.2a A1.5b B3b & B3d & A1.4a B3f & A1.4a B3b B3c & A1.4a A1.5a &  & B3.2e & B2.1a & }%
{\pfausschnitt{}{Geraden und Ebenen}{Lage und Schnittpunkt}}
\end{document}